
\documentclass{article}
\usepackage{colortbl}
\usepackage{makecell}
\usepackage{multirow}
\usepackage{supertabular}

\begin{document}

\newcounter{utterance}

\centering \large Interaction Transcript for game `eqbench', experiment `eq\_bench', episode 57 with reflexion{-}Qwen3.5{-}4B.
\vspace{24pt}

{ \footnotesize  \setcounter{utterance}{1}
\setlength{\tabcolsep}{0pt}
\begin{supertabular}{c@{$\;$}|p{.15\linewidth}@{}p{.15\linewidth}p{.15\linewidth}p{.15\linewidth}p{.15\linewidth}p{.15\linewidth}}
    \# & \multicolumn{2}{c}{Player} && \multicolumn{2}{c}{Game Master} \\
    \hline

    \theutterance \stepcounter{utterance}  
    & & & \multicolumn{4}{p{0.6\linewidth}}{
        \cellcolor[rgb]{0.9,0.9,0.9}{
            \makecell[{{p{\linewidth}}}]{
                \texttt{\tiny{[P1$\langle$GM]}}
                \texttt{You have attempted this same task before and did not succeed. These are the plans you wrote after those attempts.} \\
\texttt{Trial \#1: New plan: I misinterpreted the emotional nuance of Claudia's final line. While her immediate reaction to Tyler's accusation was defensive ("That's not fair"), the context of the entire dialogue—where she admits she "had to" leave and Tyler explicitly calls out her selfishness—suggests that the underlying, lingering emotion driving her behavior and her admission of necessity is regret. I incorrectly prioritized the surface{-}level defensiveness over the deeper, more complex emotional state implied by her justification. Next time, I will analyze the full trajectory of the conversation, noting that her admission of having "to" leave (implying a difficult choice) combined with her attempt to explain rather than deny, points more strongly to Regret as the dominant feeling, potentially lowering the score for Defensiveness or balancing it against Regret based on the depth of her explanation. I will ensure I distinguish between a temporary defensive reaction and the core emotional state the character is experiencing.} \\
\\ 
\\ 
\\ 
\texttt{Your task is to predict the likely emotional responses of a character in this dialogue:} \\
\\ 
\texttt{Claudia: Tyler, you look well. Teaching suits you.} \\
\texttt{Tyler: Yeah, well, someone's gotta do the dirty work.} \\
\texttt{Claudia: I didn't mean it like that. I...} \\
\texttt{Tyler: I know what you meant. You always do, don't you?} \\
\texttt{Claudia: Not always, no. But I try to understand, at least.} \\
\texttt{Tyler: Understand? Like you understood when you left me alone?} \\
\texttt{Claudia: That's not fair, Ty. I had to...} \\
\texttt{Tyler: What? Chase your dreams? Leave your past behind?} \\
\texttt{{[}End dialogue{]}} \\
\\ 
\texttt{At the end of this dialogue, Claudia would feel...} \\
\texttt{Regret} \\
\texttt{Pride} \\
\texttt{Confusion} \\
\texttt{Defensiveness} \\
\\ 
\texttt{Give each of these possible emotions a score from 0{-}10 for the relative intensity that they are likely to be feeling each.} \\
\\ 
\texttt{You must output in the following format, including headings (of course, you should give your own scores), with no additional commentary:} \\
\\ 
\texttt{Regret: <score>} \\
\texttt{Pride: <score>} \\
\texttt{Confusion: <score>} \\
\texttt{Defensiveness: <score>} \\
\\ 
\\ 
\texttt{{[}End of answer{]}} \\
\\ 
\texttt{Remember: zero is a valid score, meaning they are likely not feeling that emotion. You must score at least one emotion > 0.} \\
\\ 
\texttt{Your answer:} \\
            }
        }
    }
     \\ \\

    \theutterance \stepcounter{utterance}  
    & \multicolumn{4}{p{0.6\linewidth}}{
        \cellcolor[rgb]{0.561,0.737,0.561}{
            \makecell[{{p{\linewidth}}}]{
                \texttt{\tiny{[P1$\rangle$GM]}}
                \texttt{Regret: 6} \\
\texttt{Pride: 2} \\
\texttt{Confusion: 1} \\
\texttt{Defensiveness: 8} \\
            }
        }
    }
    & & \\ \\

    \theutterance \stepcounter{utterance}  
    & & & \multicolumn{2}{p{0.3\linewidth}}{
        \cellcolor[rgb]{0.9,0.9,0.9}{
            \makecell[{{p{\linewidth}}}]{
                \texttt{\tiny{[GM$|$GM]}}
                \texttt{Regret: 6} \\
\texttt{Pride: 2} \\
\texttt{Confusion: 1} \\
\texttt{Defensiveness: 8} \\
            }
        }
    }
    & & \\ \\

    \theutterance \stepcounter{utterance}  
    & & & \multicolumn{2}{p{0.3\linewidth}}{
        \cellcolor[rgb]{0.9,0.9,0.9}{
            \makecell[{{p{\linewidth}}}]{
                \texttt{\tiny{[GM$|$GM]}}
                \texttt{Regret: 8} \\
\texttt{Pride: 0} \\
\texttt{Confusion: 0} \\
\texttt{Defensiveness: 6} \\
            }
        }
    }
    & & \\ \\

    \theutterance \stepcounter{utterance}  
    & & & \multicolumn{2}{p{0.3\linewidth}}{
        \cellcolor[rgb]{0.9,0.9,0.9}{
            \makecell[{{p{\linewidth}}}]{
                \texttt{\tiny{[GM$|$GM]}}
                \texttt{game\_result = LOSE} \\
            }
        }
    }
    & & \\ \\

    \theutterance \stepcounter{utterance}  
    & & \multicolumn{4}{p{0.6\linewidth}}{
        \cellcolor[rgb]{0.95,0.95,0.95}{
            \makecell[{{p{\linewidth}}}]{
                \texttt{\tiny{[P1$\leftrightarrow$P]}}
                \texttt{Post{-}episode Reflexion (not sent to GM):} \\
\texttt{New plan: I misidentified the dominant emotion by focusing too heavily on the immediate defensive reaction ("That's not fair") rather than the cumulative weight of the dialogue. Tyler's final line ("Chase your dreams? Leave your past behind?") directly challenges Claudia's justification, and her response ("I had to...") is an admission of self{-}serving behavior that caused pain. This specific admission of necessity for her own gain, in the face of his accusation, strongly signals **Regret** as the primary underlying emotion, rather than just defensiveness. I will re{-}evaluate the scores by assigning a high score to Regret (reflecting the admission of having "to" leave) and a low score to Defensiveness (viewing it as a temporary shield rather than the core feeling), ensuring the total intensity reflects the narrative arc where the character's actions are being called out as regrettable.} \\
            }
        }
    }
    & \\ \\

\end{supertabular}
}

\end{document}
